% Figure: a Hohmann transfer between two coplanar circular orbits, with
% the two impulsive burns at the tangent points and the transfer ellipse
% touching the inner orbit at periapsis and the outer orbit at apoapsis.
\begin{figure}[htbp]
\centering
\begin{tikzpicture}[
  burn/.style={-{Stealth[length=3mm]}, very thick, draw=engred},
  lbl/.style={font=\small},
]
% central body
\fill[engamber] (0,0) circle (5pt);
\node[engamber, lbl, below left] at (0,0) {Earth};
% inner circular orbit, radius 1.7
\draw[engblue, very thick] (0,0) circle (1.7);
\node[engblue, lbl] at (-1.7,-1.35) {inner orbit $r_1$};
% outer circular orbit, radius 3.4
\draw[enggray, very thick] (0,0) circle (3.4);
\node[enggray, lbl] at (2.4,2.7) {outer orbit $r_2$};
% transfer ellipse: periapsis at inner orbit (left, x=-1.7),
% apoapsis at outer orbit (right, x=+3.4).
% semi-major a = (1.7+3.4)/2 = 2.55; centre offset along x:
% centre x = apoapsis - a = 3.4 - 2.55 = 0.85.
% b = sqrt(r1*r2) = sqrt(1.7*3.4) = sqrt(5.78) = 2.404.
\draw[englime!70!black, very thick, dashed] (0.85,0) ellipse (2.55 and 2.404);
\node[englime!60!black, lbl] at (0.85,-2.7) {transfer ellipse};
% burn 1 at periapsis (left tangent point on inner orbit), x=-1.7,y=0
\coordinate (b1) at (-1.7,0);
\fill[engblack] (b1) circle (1.6pt);
\draw[burn] (b1) ++(0,0) -- ++(0,0.9);
\node[engred, lbl, left] at (-1.7,0.5) {$\Delta v_1$};
% burn 2 at apoapsis (right tangent point on outer orbit), x=+3.4,y=0
\coordinate (b2) at (3.4,0);
\fill[engblack] (b2) circle (1.6pt);
\draw[burn] (b2) ++(0,0) -- ++(0,0.9);
\node[engred, lbl, right] at (3.4,0.5) {$\Delta v_2$};
% direction-of-motion arrows
\draw[-{Stealth[length=2.5mm]}, engblue, thick] (0,-1.7) ++(0.25,0) arc (0:-22:1.7);
\draw[-{Stealth[length=2.5mm]}, enggray, thick] (0,3.4) ++(-0.25,0) arc (90:112:3.4);
\end{tikzpicture}
\caption[Hohmann transfer between two circular orbits]{A Hohmann
transfer raises a spacecraft from an inner circular orbit of radius
$r_1$ to an outer circular orbit of radius $r_2$ with two prograde
burns. The first burn $\Delta v_1$ at periapsis injects the craft onto
an ellipse whose periapsis touches the inner orbit and whose apoapsis
touches the outer orbit; the craft coasts half an ellipse to apoapsis,
where the second burn $\Delta v_2$ circularises onto the outer orbit.
The transfer ellipse has semi-major axis $a=(r_1+r_2)/2$, and the
coast lasts half its period. The Hohmann transfer is the
minimum-energy two-impulse transfer between coplanar circular orbits,
which is why it sets the reference $\Delta v$ against which other
transfer schemes are measured.}
\label{fig:vol03ch04:hohmann}
\end{figure}
